\section*{Section 1 Concept Quiz}
\addcontentsline{toc}{section}{Section 1 Concept Quiz}

Use this as a short retrieval check after finishing the section. On the website, each item should accept an answer before revealing the hint and worked explanation.

\begin{bgcheck}{limit-meaning-q1}{integer}{7}[Concept check]
If \(\lim_{x\to 4}f(x)=7\), what output value is approached?
\begin{bghint}
The number after the equals sign is the output target.
\end{bghint}
\begin{bgreveal}
\(7\). The input approaches \(4\), while the output approaches \(7\).
\end{bgreveal}
\end{bgcheck}

\begin{bgcheck}{one-sided-q1}{integer}{3}[Concept check]
A graph approaches \(3\) from the left at \(x=2\). What is \(\lim_{x\to2^-}f(x)\)?
\begin{bghint}
The superscript minus means use inputs smaller than \(2\).
\end{bghint}
\begin{bgreveal}
The left-hand limit is \(3\).
\end{bgreveal}
\end{bgcheck}

\begin{bgcheck}{hole-q1}{integer}{5}[Concept check]
A curve approaches the open point \((1,5)\), while a filled point is at \((1,9)\). What is the limit as \(x\to1\)?
\begin{bghint}
The limit follows the nearby curve, not the filled point.
\end{bghint}
\begin{bgreveal}
The limit is \(5\); the function value is \(9\).
\end{bgreveal}
\end{bgcheck}

\begin{bgcheck}{rate-q1}{integer}{12}[Concept check]
An average-rate expression simplifies to \(12-2h\). What value does it approach as \(h\to0\)?
\begin{bghint}
Set only the limiting value of \(h\) after simplification.
\end{bghint}
\begin{bgreveal}
\(12-2(0)=12\).
\end{bgreveal}
\end{bgcheck}

\begin{bgsummary}[title={Quiz use on the website}]
This page should report completion by concept, not merely a total score. A wrong answer should link back to the exact lesson that teaches the missing method. The same questions may appear in randomized numerical variants, but the canonical explanation should remain stable.
\end{bgsummary}
